\chapter{Glosario de términos y acrónimos}
\label{anexo:glosario_acronimos}

Este anexo recoge los principales términos técnicos y acrónimos utilizados en la memoria de TierraViva. Su objetivo es facilitar la lectura del documento, especialmente en los apartados donde confluyen conceptos de electrónica, comunicaciones móviles, sistemas embebidos, plataformas IoT, bases de datos, API y desarrollo web.

El glosario se complementa con los anexos técnicos específicos. La estructura de campos de ThingSpeak se documenta en el Anexo \ref{anexo:estructura_thingspeak_estados}; la referencia de \textit{endpoints}, en el Anexo \ref{anexo:manual_api}; y el uso del \textit{dashboard}, en el Anexo \ref{anexo:manual_usuario}.

\section{Glosario de términos}

\begin{description}

    \item[A7670E-LASA]
    Módulo de comunicaciones móviles LTE utilizado por la estación remota para publicar telemetría en ThingSpeak mediante peticiones HTTP.

    \item[Actuador]
    Elemento capaz de ejecutar una acción física sobre el entorno. En TierraViva, la bomba de agua de 12 V actúa como actuador del sistema de riego.

    \item[Adafruit IO]
    Plataforma IoT orientada a proyectos de electrónica conectada. Se consideró como alternativa para publicación y consulta de datos, aunque finalmente se seleccionó ThingSpeak por su sencillez de integración con la arquitectura propuesta.

    \item[Agricultura de precisión]
    Enfoque agrícola basado en medir variables del entorno y del cultivo para tomar decisiones más ajustadas que un riego fijo o manual.

    \item[Agrolab Móstoles]
    Entorno de agricultura urbana/periurbana situado en Móstoles, Madrid, utilizado en TierraViva para realizar una prueba funcional del prototipo fuera del entorno de laboratorio.

    \item[ALTER TABLE]
    Instrucción SQL utilizada para modificar la estructura de una tabla existente. En TierraViva se emplea en la estrategia de migración suave del esquema SQLite durante la evolución del prototipo.

    \item[API]
    Interfaz que permite consultar o intercambiar datos con una aplicación. En TierraViva, la API FastAPI expone configuración, última lectura, histórico, estadísticas y estado del sistema.

    \item[API Key]
    Clave utilizada para leer o escribir datos en un servicio. En TierraViva se emplean claves de ThingSpeak, que se mantienen fuera de la memoria por motivos de seguridad.

    \item[AQUASTAT]
    Base de datos internacional sobre agua y agricultura utilizada como fuente para contextualizar el peso del uso agrícola en las extracciones mundiales de agua.
    
    \item[Arduino UNO]
    Placa de microcontrolador utilizada como unidad de control de la estación remota. Lee sensores, gestiona el módem LTE y controla el relé.

    \item[ATmega328P]
    Microcontrolador principal de la placa Arduino UNO compatible utilizada en el prototipo. Ejecuta el \textit{firmware} de estación y coordina sensores, módem LTE, relé y lógica local de riego.

    \item[Automatización basada en reglas]
    Estrategia de decisión en la que una acción se ejecuta si se cumplen condiciones explícitas. En TierraViva, el riego se decide mediante umbrales, validación de lecturas y condiciones de seguridad.

    \item[Backend]
    Parte del software que ejecuta la lógica interna del sistema. En TierraViva se ejecuta en la Raspberry Pi e incluye la ingesta de datos, SQLite y la API.

    \item[Batería 12 V]
    Fuente de alimentación principal utilizada para proporcionar energía a la estación y a la carga de actuación. En TierraViva alimenta el sistema de 12 V y sirve de entrada al regulador LM2596.

    \item[Bash]
    Intérprete de comandos utilizado en sistemas Linux para ejecutar instrucciones y automatizar tareas. En TierraViva se emplea en el script \texttt{tierraviva.sh} para instalación, arranque, parada, estado y consulta de logs.

    \item[Baudio]
    Unidad asociada a la velocidad de señalización en una comunicación serie. En TierraViva se utiliza una comunicación a 9600 baudios entre Arduino y el módem para mejorar la estabilidad con SoftwareSerial.

    \item[BH1750]
    Sensor digital de luminosidad utilizado para medir luz ambiental en lux. En algunos módulos aparece integrado como GY-302.

    \item[Blynk]
    Plataforma IoT orientada a la creación de aplicaciones y paneles de control para dispositivos conectados. Se valoró como alternativa de visualización y gestión, pero no se adoptó en la versión final para mantener una arquitectura más controlada y documentable desde Raspberry Pi.

    \item[BME280]
    Sensor ambiental utilizado para medir temperatura, humedad relativa y presión atmosférica mediante bus I2C.

    \item[Bomba de agua]
    Actuador hidráulico utilizado para mover agua desde un depósito hacia el cultivo o sustrato durante una prueba de riego.

    \item[Broker]
    Elemento funcional intermedio que desacopla productores y consumidores de información. Recibe datos o mensajes procedentes de unos componentes del sistema y los pone a disposición de otros mediante mecanismos de publicación, consulta o suscripción. Aunque el término se utiliza con frecuencia en arquitecturas MQTT, también puede aplicarse de forma funcional a otros mecanismos de intercambio de datos. En TierraViva, ThingSpeak actúa como \textit{broker} de telemetría entre las estaciones de campo y la Raspberry Pi.

    \item[Búfer]
    Zona de memoria utilizada para almacenar temporalmente datos o cadenas. En Arduino UNO debe dimensionarse con cuidado por la memoria SRAM limitada. 

    \item[Bus I2C]
    Bus de comunicación serie empleado para conectar sensores digitales con un microcontrolador. En TierraViva lo utilizan el BME280 y el BH1750.

    \item[C/C++]
    Lenguajes utilizados en el entorno Arduino para desarrollar el \textit{firmware} de las estaciones. En TierraViva se emplean para implementar la lectura de sensores, la comunicación con el módem LTE y la lógica local de riego.

    \item[Cable Dupont]
    Cable de conexión utilizado habitualmente en prototipos electrónicos para unir pines de sensores, placas de desarrollo, protoboards y módulos auxiliares.

    \item[Campaña agronómica]
    Evaluación prolongada de un sistema en condiciones reales de cultivo durante un periodo suficientemente amplio para estudiar comportamiento estacional, respuesta del cultivo y ajuste fino de riego. En TierraViva, la prueba realizada se documenta como validación funcional del prototipo en entorno real, dejando una campaña agronómica completa como posible fase futura.

    \item[Canal descendente]
    Canal de comunicación desde el sistema de supervisión hacia la estación. En TierraViva se plantea como posible evolución futura para enviar configuración u órdenes supervisadas, manteniendo en la versión actual la decisión de riego dentro del \textit{firmware} de cada estación.

    \item[Canal ThingSpeak]
    Espacio de almacenamiento en ThingSpeak donde una estación publica sus lecturas y estados. Cada estación debe tener su propio canal para evitar mezclar datos.

    \item[Caso de prueba]
    Comprobación definida para validar una funcionalidad, condición o comportamiento concreto del sistema. En TierraViva los casos de prueba se identifican mediante códigos P01--P17.

    \item[Caudal]
    Volumen de agua que circula por unidad de tiempo. En TierraViva se plantea como posible variable futura para cuantificar el volumen real aplicado durante cada riego.

    \item[Credenciales]
    Datos de autenticación, como claves de lectura o escritura, utilizados para acceder a servicios o canales de comunicación.

    \item[CH340G]
    Conversor USB-serie presente en la placa Arduino UNO compatible utilizada. Permite la comunicación entre el ordenador y el microcontrolador para programar y depurar el \textit{firmware}.

    \item[Cirkit Designer]
    Herramienta de diseño y documentación de circuitos utilizada para elaborar el esquema eléctrico general de una estación TierraViva.

    \item[Clave primaria]
    Campo o conjunto de campos que identifica de forma única cada registro de una tabla. En el modelo SQLite de TierraViva se representa mediante la abreviatura PK.

    \item[Cloud]
    Infraestructura o servicio remoto usado para almacenar, procesar o intercambiar datos. En TierraViva, ThingSpeak se usa como plataforma de telemetría y consulta entre las estaciones de campo y la Raspberry Pi.

    \item[CoAP]
    Protocolo de aplicación orientado a dispositivos y redes con recursos limitados. Se estudia como alternativa, pero no se utiliza en la implementación actual.

    \item[Comando AT]
    Instrucción enviada a un módem para configurarlo o solicitar información. En TierraViva se usan comandos AT para gestionar el A7670E-LASA.

    \item[Condensador]
    Componente electrónico utilizado para almacenar carga eléctrica de forma temporal y ayudar a estabilizar tensiones o filtrar variaciones rápidas. En TierraViva se menciona como elemento auxiliar de estabilización eléctrica.

    \item[Conductividad eléctrica]
    Magnitud relacionada con la concentración de sales disueltas en el agua o el sustrato. Se menciona como posible ampliación futura para mejorar la caracterización agronómica del entorno.

    \item[Contacto COM]
    Contacto común del relé. En el bloque de actuación se conecta con el contacto normalmente abierto cuando el relé se activa, cerrando el circuito de alimentación de la bomba.

    \item[Control descendente]
    Envío de órdenes desde una plataforma o \textit{dashboard} hacia una estación remota. En TierraViva queda planteado como línea futura, condicionada al diseño de validación de órdenes, caducidad de comandos, límites de seguridad y confirmación de ejecución.

    \item[Controlador CH340/CH341]
    Controlador necesario para que Windows reconozca correctamente algunas placas Arduino compatibles basadas en conversores USB-serie CH340 o CH341.

    \item[Cooldown]
    Periodo mínimo de seguridad entre dos activaciones consecutivas del riego. Evita que la bomba se active repetidamente en intervalos demasiado cortos.

    \item[curl]
    Herramienta de línea de comandos utilizada para realizar peticiones HTTP. En TierraViva se emplea como forma sencilla de comprobar respuestas de la API local, por ejemplo en los endpoints \texttt{/api/config}, \texttt{/api/latest} o \texttt{/api/health}.

    \item[\textit{Dashboard}]
    Interfaz visual del sistema. En TierraViva muestra lecturas actuales, gráficas, estado de conexión, eventos, diagnóstico e histórico reciente.

    \item[Dato obsoleto]
    Lectura almacenada cuya antigüedad supera el margen configurado. Puede indicar que la estación no está publicando datos recientemente o que existe un problema de ingesta.

    \item[Decisión local]
    Lógica ejecutada directamente por la estación Arduino. En TierraViva, la decisión de riego se toma en la estación a partir de las lecturas de sensores, los umbrales configurados y las reglas de seguridad implementadas en el \textit{firmware}.

    \item[Divisor resistivo]
    Circuito formado por resistencias que permite reducir una tensión. En TierraViva se utiliza para adaptar la señal serie enviada desde Arduino hacia la entrada RXD del módulo A7670E-LASA.

    \item[Endpoint]
    Punto de acceso expuesto por una API para realizar una operación concreta mediante HTTP. En TierraViva, estos puntos de acceso permiten consultar configuración, última lectura, histórico, estadísticas y estado del sistema.

    \item[Estado seguro]
    Condición conservadora del sistema en la que la actuación física se mantiene desactivada o limitada para evitar riesgos. En TierraViva implica relé apagado, bloqueo de riegos repetidos y bomba activa solo durante intervalos temporizados.

    \item[Entorno virtual]
    Entorno aislado de Python utilizado para instalar dependencias del proyecto sin modificar la instalación global del sistema operativo.

    \item[Entry ID]
    Identificador de entrada asignado por ThingSpeak a cada publicación. La Raspberry Pi lo utiliza para evitar almacenar lecturas duplicadas.

    \item[ESP32]
    Familia de microcontroladores considerada como alternativa a Arduino UNO por su mayor potencia de cálculo y conectividad integrada.

    \item[Estación remota]
    Nodo físico instalado en campo o entorno de pruebas. En TierraViva existen dos estaciones remotas implementadas, A y B, formadas por Arduino UNO, sensores, módulo LTE, relé, alimentación y sistema de riego.

    \item[Evapotranspiración]
    Pérdida de agua asociada conjuntamente a la evaporación del suelo y a la transpiración de las plantas. En TierraViva se menciona como posible variable futura para ajustar mejor el riego a las condiciones meteorológicas y del cultivo.

    \item[FastAPI]
    Framework de Python utilizado para crear la API local de TierraViva.

    \item[Fichero .env]
    Archivo de configuración utilizado para cargar variables de entorno durante la ejecución del sistema. Permite mantener claves y parámetros fuera del código fuente.

    \item[\textit{Firmware}]
    Programa cargado en el microcontrolador de la estación. Controla sensores, comunicación LTE, publicación de datos, relé y lógica local de riego.

    \item[Frontend]
    Parte visual de una aplicación web. En TierraViva corresponde al \textit{dashboard} desarrollado con HTML, CSS y JavaScript.

    \item[Frescura de datos]
    Clasificación de la antigüedad de la última lectura. La API y el \textit{dashboard} la utilizan para diferenciar datos recientes, antiguos u offline.

    \item[Gateway]
    Equipo que actúa como pasarela entre dispositivos y una red externa. Algunas tecnologías como LoRaWAN pueden requerir infraestructura de gateway.

    \item[GET]
    Método HTTP utilizado para solicitar o leer un recurso desde un servidor. En TierraViva se menciona como una de las formas habituales de interacción con interfaces HTTP.

    \item[Git]
    Sistema de control de versiones utilizado para registrar cambios en el código fuente y facilitar la organización del desarrollo.

    \item[GitHub]
    Plataforma de alojamiento de repositorios Git utilizada para conservar el código fuente, documentación técnica y evidencias del proyecto TierraViva.

    \item[GY-302]
    Módulo comercial basado en el sensor BH1750 para medir luminosidad ambiente.

    \item[Hardware]
    Conjunto de componentes físicos del sistema. En TierraViva incluye estaciones remotas, sensores, módulo LTE, relé, bomba, alimentación, cableado y Raspberry Pi.

    \item[hPa]
    Hectopascal, unidad de presión utilizada para expresar la presión atmosférica medida por el sensor BME280. En TierraViva aparece asociada a la variable \texttt{pressure\_hpa}.

    \item[Ingesta de datos]
    Proceso mediante el cual la Raspberry Pi recupera lecturas desde ThingSpeak y las almacena en SQLite.

    \item[Ingestor]
    Componente software encargado de recuperar lecturas desde ThingSpeak, normalizarlas y almacenarlas en SQLite. En TierraViva corresponde principalmente al proceso implementado en \texttt{main.py}.

    \item[IoT]
    Paradigma tecnológico basado en conectar objetos físicos a Internet para medir, comunicar, procesar o actuar sobre el entorno.

    \item[ISO 8601]
    Estándar internacional para representar fechas y horas de forma textual. En TierraViva se utiliza en marcas temporales devueltas por la API, como \texttt{ts}, \texttt{first\_ts} o \texttt{last\_ts}.

    \item[JavaScript]
    Lenguaje utilizado en el frontend del dashboard para consultar la API, procesar respuestas y actualizar dinámicamente la interfaz web.

    \item[Latencia]
    Tiempo de retardo entre el envío de una orden, dato o petición y la obtención de una respuesta o efecto observable. En TierraViva se menciona al analizar el control remoto descendente de riego, ya que una actuación sobre agua requiere tiempos controlados, trazabilidad y mecanismos de seguridad.

    \item[LaTeX]
    Sistema de composición de documentos utilizado para redactar la memoria, anexos, tablas, figuras, referencias y estructura formal del TFG.

    \item[Lectura]
    Registro asociado a una estación en un instante concreto. Puede incluir humedad del suelo, temperatura, humedad ambiental, luminosidad, presión, estado del relé, evento y diagnóstico.

    \item[LM2596/LM2596S]
    Regulador DC-DC de tipo reductor utilizado para convertir una tensión de entrada superior, como 12 V, en una salida regulada de 5 V.

    \item[Log]
    Registro de ejecución generado por un programa o servicio. En TierraViva los logs permiten revisar errores de ingesta, arranque, configuración o ejecución de la API.

    \item[LoRa]
    Modulación radio de largo alcance y bajo consumo considerada como alternativa de comunicación para estaciones remotas. En TierraViva se analiza junto con LoRaWAN, aunque finalmente se adopta conectividad LTE para validar el prototipo completo dentro del alcance del TFG.

    \item[LoRaWAN]
    Tecnología LPWAN considerada como alternativa de comunicación. En TierraViva se prioriza LTE para la versión desarrollada por su disponibilidad inmediata mediante SIM, integración con el módulo A7670E-LASA y compatibilidad con publicación HTTP hacia ThingSpeak.

    \item[LPWAN]
    Familia de redes inalámbricas de bajo consumo y largo alcance utilizadas en IoT. Incluye tecnologías como LoRaWAN, NB-IoT y LTE-M.

    \item[LTE]
    Tecnología de comunicación móvil utilizada por el módulo A7670E-LASA para proporcionar conectividad a la estación remota.

    \item[Macarrón termorretráctil]
    Material aislante que se contrae al aplicar calor y se utiliza para proteger uniones eléctricas, cables y soldaduras. En TierraViva se emplea como material auxiliar de aislamiento y protección mecánica.

    \item[Masa común]
    Referencia eléctrica compartida entre los distintos módulos del circuito. En TierraViva permite que Arduino, sensores, relé, regulador y módulo LTE interpreten correctamente las señales entre sí.

    \item[Modo seguro]
    Estado conservador de funcionamiento en el que la estación evita actuaciones peligrosas o no justificadas ante lecturas no fiables o condiciones anómalas.

    \item[Monitor serie]
    Herramienta del entorno Arduino utilizada para visualizar mensajes enviados por la placa durante la ejecución del \textit{firmware}. En TierraViva permite comprobar lecturas de sensores, respuestas del módem, registro LTE y publicación en ThingSpeak.

    \item[Mosquitto]
    \textit{Broker} MQTT utilizado en el prototipo local previo para validar la cadena de adquisición, transporte, almacenamiento y visualización antes de adoptar la arquitectura final basada en LTE y ThingSpeak. En ese prototipo, Mosquitto actuaba como intermediario de mensajes bajo el modelo publicación/suscripción propio de MQTT.

    \item[MQTT]
    Protocolo de mensajería publicación/suscripción habitual en IoT. Se estudia como alternativa, pero TierraViva utiliza HTTP en la versión actual.

    \item[Módulo relé]
    Dispositivo que permite controlar una carga eléctrica mediante una señal de baja potencia. En TierraViva se utiliza para activar o desactivar la bomba de agua.
    
    \item[Multímetro]
    Instrumento de medida utilizado para comprobar tensiones, continuidad y referencias eléctricas durante el montaje. En TierraViva se emplea especialmente para validar alimentación, regulador, masa común y conexiones antes de integrar los módulos.

    \item[Nivel de depósito]
    Medida asociada a la cantidad de agua disponible en el recipiente o sistema de almacenamiento. Se plantea como mejora futura para detectar falta de agua antes de activar la bomba.

    \item[Nodo de campo]
    Dispositivo situado en el entorno físico de medida. En TierraViva equivale a la estación remota basada en Arduino UNO.

    \item[Nodo de supervisión]
    Equipo encargado de recoger, almacenar y mostrar información del sistema. En TierraViva este papel lo cumple la Raspberry Pi.

    \item[Optoacoplador]
    Componente que permite transmitir una señal entre dos partes de un circuito mediante aislamiento óptico. En TierraViva aparece integrado en el módulo relé para mejorar la separación entre la señal de control y la carga conmutada.

    \item[pH]
    Magnitud utilizada para expresar el carácter ácido o básico de una disolución. En TierraViva se menciona como posible variable futura para ampliar la caracterización agronómica del entorno.

    \item[Plataforma cloud]
    Servicio remoto que permite recibir, almacenar o consultar datos. En TierraViva, ThingSpeak actúa como plataforma cloud de telemetría.

    \item[POST]
    Método HTTP utilizado para enviar datos a un servidor. En TierraViva se menciona en el contexto de publicación de telemetría o envío de información hacia plataformas compatibles con HTTP.

    \item[Protoboard]
    Placa de pruebas sin soldadura utilizada para realizar montajes temporales y validar conexiones durante el desarrollo del prototipo.

    \item[Prototipo académico]
    Versión funcional desarrollada dentro del alcance de un Trabajo de Fin de Grado. En TierraViva permite validar arquitectura, sensorización, comunicación, almacenamiento, visualización y actuación como base técnica para futuras versiones más robustas.

    \item[Prueba funcional en entorno real]
    Ensayo realizado fuera del montaje de laboratorio para comprobar que el prototipo mantiene su flujo principal de funcionamiento en un contexto de uso más próximo al previsto. En TierraViva se aplica a la prueba realizada en Agrolab Móstoles.

    \item[Publicación de telemetría]
    Envío de lecturas desde la estación hacia una plataforma externa. En TierraViva se realiza mediante HTTP sobre LTE hacia ThingSpeak.

    \item[Puente serie--MQTT]
    Componente software utilizado en el prototipo local previo para adaptar lecturas recibidas por puerto serie y publicarlas como mensajes MQTT. En la arquitectura final se sustituye por publicación LTE/HTTP hacia ThingSpeak.

    \item[Puerto COM]
    Identificador de puerto serie utilizado por Windows para comunicarse con dispositivos conectados por USB-serie. En TierraViva se emplea para programar y depurar Arduino desde Arduino IDE.

    \item[Python]
    Lenguaje utilizado en la Raspberry Pi para implementar la ingesta de datos, el acceso a SQLite y la API FastAPI del sistema TierraViva.

    \item[Raspberry Pi]
    Ordenador de placa reducida utilizado como nodo local de supervisión, almacenamiento, API y \textit{dashboard}.

    \item[Raw JSON]
    Respuesta original de ThingSpeak conservada en la base de datos. Permite mantener trazabilidad de los campos recibidos, incluso si no todos se almacenan como columnas explícitas.

    \item[Registro histórico]
    Conjunto de lecturas almacenadas a lo largo del tiempo. Permite consultar evolución, estadísticas y comportamiento del sistema.

    \item[Relé normalmente abierto]
    Configuración del relé en la que el circuito permanece abierto cuando no hay señal de activación. Es la opción más segura para mantener el riego apagado por defecto.
    
    \item[Repositorio]
    Espacio donde se almacena el código fuente y la documentación técnica del proyecto. En TierraViva contiene el \textit{firmware}, el software de Raspberry Pi, la API, el \textit{dashboard} y los ficheros auxiliares de instalación.

    \item[Resistencia]
    Componente electrónico que limita o ajusta el paso de corriente. En TierraViva se utiliza especialmente en el divisor resistivo empleado para adaptar niveles de señal.

    \item[Restricción UNIQUE]
    Restricción de base de datos que impide insertar registros duplicados para una combinación de campos. En TierraViva se aplica sobre \texttt{station} y \texttt{entry\_id} para evitar duplicar lecturas ya procesadas.

    \item[Riego inteligente]
    Riego que utiliza medidas del entorno o del suelo para decidir cuándo actuar, en lugar de basarse únicamente en horarios fijos.

    \item[Riesgo residual]
    Riesgo que permanece tras aplicar medidas de mitigación. En TierraViva se utiliza para delimitar los aspectos que siguen presentes en el prototipo académico y para orientar mejoras futuras en robustez, seguridad, autonomía y validación prolongada.

    \item[Script]
    Programa auxiliar utilizado para automatizar tareas de ejecución, instalación o mantenimiento. En TierraViva se utiliza especialmente en el fichero \texttt{tierraviva.sh}.

    \item[Seguridad funcional]
    Conjunto de medidas orientadas a evitar actuaciones peligrosas o no deseadas del sistema. En TierraViva se aplica especialmente al control del relé, la bomba, los tiempos máximos de riego, el modo seguro y el bloqueo por \textit{cooldown}.

    \item[Sensor capacitivo de humedad del suelo]
    Sensor que permite estimar la humedad del sustrato mediante una lectura analógica. Es la variable principal para la decisión local de riego.

    \item[Sensorización]
    Conjunto de sensores y medidas utilizadas para observar el estado del entorno físico.

    \item[Sketch]
    Nombre habitual de un programa desarrollado para Arduino. En TierraViva corresponde al \textit{firmware} cargado en cada estación.

    \item[Software]
    Conjunto de programas y ficheros ejecutables del sistema. En TierraViva incluye el \textit{firmware} de estación, el ingestor, la base de datos, la API, el \textit{dashboard} y los scripts de gestión.

    \item[SoftwareSerial]
    Librería de Arduino utilizada para crear un puerto serie software en pines digitales. En TierraViva permite comunicar Arduino UNO con el módulo A7670E-LASA sin ocupar el puerto serie hardware de programación.

    \item[SQLite]
    Sistema de base de datos ligero basado en fichero. En TierraViva almacena lecturas, estados y eventos de forma local en la Raspberry Pi.

    \item[STM32]
    Familia de microcontroladores considerada como alternativa por sus mayores prestaciones y flexibilidad, aunque con una complejidad de configuración superior para el alcance del TFG.

    \item[Tarjeta microSD]
    Tarjeta de memoria de pequeño formato utilizada como almacenamiento principal en la Raspberry Pi. En TierraViva se emplea para alojar el sistema operativo, el software del proyecto, la base de datos SQLite y los ficheros auxiliares de ejecución.

    \item[Telemetría]
    Conjunto de datos medidos por una estación y enviados a un sistema de supervisión.

    \item[ThingsBoard]
    Plataforma IoT de código abierto orientada a la gestión de dispositivos, telemetría, reglas y \textit{dashboards}. Se consideró como alternativa más completa, aunque se descartó para el alcance actual por requerir mayor configuración e infraestructura.

    \item[ThingSpeak]
    Plataforma IoT de MathWorks utilizada en TierraViva como \textit{broker} de telemetría y estado. Recibe los datos publicados por las estaciones de campo y los pone a disposición de la Raspberry Pi para su consulta, almacenamiento local y visualización.
    
    \item[Topic]
    Tema o canal lógico utilizado en MQTT para organizar mensajes publicados y suscripciones.

    \item[Trazabilidad]
    Capacidad de reconstruir el origen y evolución de un dato, decisión o evento. En TierraViva se apoya en \texttt{entry\_id}, \texttt{raw\_json}, SQLite, logs y documentación de pruebas.

    \item[Ubidots]
    Plataforma IoT orientada a la gestión, visualización y análisis de datos de dispositivos conectados. Se valoró como alternativa \textit{cloud}; en TierraViva se priorizó ThingSpeak por su sencillez, verificabilidad y compatibilidad con el flujo HTTP planteado.

    \item[Umbral]
    Valor límite utilizado para tomar una decisión. En TierraViva se usan umbrales de humedad del suelo para clasificar el estado del sustrato y decidir si puede ser necesario regar.

    \item[USB-C]
    Tipo de conector USB reversible utilizado por algunos dispositivos modernos. En TierraViva aparece asociado a la alimentación de la Raspberry Pi.

    \item[Uvicorn]
    Servidor ASGI utilizado para ejecutar la aplicación FastAPI de TierraViva.

    \item[Validación incremental]
    Estrategia de pruebas por bloques. Permite comprobar sensores, comunicaciones, almacenamiento, API, \textit{dashboard} y actuación antes de validar el sistema completo.

    \item[Variable de entorno]
    Parámetro de configuración definido fuera del código fuente. En TierraViva se utiliza para almacenar identificadores de canal, claves de lectura y otros valores necesarios para la ejecución en Raspberry Pi.

\end{description}

\section{Acrónimos}

\renewcommand{\arraystretch}{1.45}
\scriptsize

\begin{longtable}{|p{0.14\textwidth}|p{0.30\textwidth}|p{0.46\textwidth}|}
\hline
\rowcolor{URJCRed}
\TVHead{Acrónimo} & \TVHead{Significado} & \TVHead{Uso en TierraViva} \\
\hline
\endfirsthead

\hline
\rowcolor{URJCRed}
\TVHead{Acrónimo} & \TVHead{Significado} & \TVHead{Uso en TierraViva} \\
\hline
\endhead

\hline
\multicolumn{3}{r}{\small Continúa en la página siguiente} \\
\endfoot

\hline
\endlastfoot

ADC & \textit{Analog-to-Digital Converter} & Conversión de la lectura analógica del sensor de humedad del suelo. \\
\hline
API & \textit{Application Programming Interface} & Interfaz HTTP expuesta por FastAPI. \\
\hline
APN & \textit{Access Point Name} & Parámetro necesario para que el módulo LTE acceda a Internet mediante la SIM. \\
\hline
ASGI & \textit{Asynchronous Server Gateway Interface} & Interfaz utilizada por Uvicorn para ejecutar FastAPI. \\
\hline
ASIN & \textit{Amazon Standard Identification Number} & Identificador comercial utilizado por Amazon para referenciar productos. En TierraViva se emplea para documentar la trazabilidad de algunos componentes adquiridos. \\
\hline
AT & \textit{Attention Command} & Comandos enviados al módem A7670E-LASA. \\
\hline
BLE & \textit{Bluetooth Low Energy} & Tecnología inalámbrica de bajo consumo estudiada en el contexto general de comunicaciones IoT. \\
\hline
CoAP & \textit{Constrained Application Protocol} & Protocolo IoT ligero considerado en el estado del arte como alternativa a HTTP y MQTT. \\
\hline
CSS & \textit{Cascading Style Sheets} & Hojas de estilo del \textit{dashboard} web. \\
\hline
CSV & \textit{Comma-Separated Values} & Formato tabular útil para exportación o tratamiento de datos, aunque no es el formato principal del sistema actual. \\
\hline
DAFO & Debilidades, Amenazas, Fortalezas y Oportunidades & Herramienta de análisis estratégico utilizada para valorar factores internos y externos del proyecto TierraViva. \\
\hline
DC & \textit{Direct Current} & Corriente continua usada en la alimentación de 12 V y 5 V del prototipo. \\
\hline
ET & Especificación temporal y operativa & Identificador utilizado para los parámetros temporales y criterios operativos del sistema, como frecuencias de lectura, envío de telemetría, tiempo máximo de bomba y estado seguro. \\
\hline
GB & \textit{Gigabyte} & Unidad de capacidad utilizada para indicar la memoria de la Raspberry Pi 5 empleada en el prototipo. \\
\hline
GND & \textit{Ground} & Referencia eléctrica común del circuito. \\
\hline
GPIO & \textit{General Purpose Input/Output} & Pines de entrada/salida usados para control digital. \\
\hline
HTML & \textit{HyperText Markup Language} & Estructura del \textit{dashboard} web. \\
\hline
HTTP & \textit{HyperText Transfer Protocol} & Protocolo usado para publicar datos en ThingSpeak y consultar la API local. \\
\hline
IA & Inteligencia artificial & Línea planteada como posible evolución futura, vinculada a la disponibilidad de histórico suficiente y validación prolongada. \\
\hline
I2C & \textit{Inter-Integrated Circuit} & Bus usado por los sensores BME280 y BH1750. \\
\hline
ID & Identificador & Código utilizado para distinguir estaciones, canales, entradas, registros o parámetros dentro de la documentación y configuración de TierraViva. \\
\hline
IDE & \textit{Integrated Development Environment} & Entorno utilizado para programar y cargar el \textit{firmware} en Arduino. \\
\hline
IEEE & \textit{Institute of Electrical and Electronics Engineers} & Organización profesional y editorial citada a través de \textit{IEEE Sensors Journal} en el estado del arte sobre riego inteligente e IoT. \\
\hline
IoT & \textit{Internet of Things} & Enfoque general del sistema TierraViva. \\
\hline
IP & \textit{Internet Protocol} & Dirección utilizada para acceder al \textit{dashboard} o a la API de la Raspberry Pi. \\
\hline
JSON & \textit{JavaScript Object Notation} & Formato de datos utilizado por la API y por las respuestas de ThingSpeak. \\
\hline
LED & \textit{Light Emitting Diode} & Indicador luminoso empleado en el prototipo. \\
\hline
LoRaWAN & \textit{Long Range Wide Area Network} & Tecnología LPWAN considerada como alternativa de comunicación. \\
\hline
LPWAN & \textit{Low-Power Wide-Area Network} & Familia de redes IoT de bajo consumo y largo alcance. \\
\hline
LTE & \textit{Long Term Evolution} & Tecnología móvil utilizada por el módulo A7670E-LASA. \\
\hline
LTE-M & \textit{Long Term Evolution for Machines} & Tecnología celular orientada a comunicaciones IoT con mayor movilidad y menor latencia que NB-IoT. \\
\hline
M2M & \textit{Machine to Machine} & Comunicación entre dispositivos, relevante en protocolos como MQTT. \\
\hline
MQTT & \textit{Message Queuing Telemetry Transport} & Protocolo publicación/suscripción estudiado como alternativa a HTTP. \\
\hline
NB-IoT & \textit{Narrowband Internet of Things} & Tecnología celular LPWAN considerada dentro del contexto de comunicaciones IoT. \\
\hline
NO & \textit{Normally Open} & Contacto normalmente abierto del relé. En TierraViva se utiliza para que la bomba permanezca apagada mientras no exista una activación expresa desde el \textit{firmware}. \\
\hline
OASIS & \textit{Organization for the Advancement of Structured Information Standards} & Organización de estandarización responsable de la especificación MQTT. \\
\hline
ODS & Objetivos de Desarrollo Sostenible & Marco de Naciones Unidas utilizado para relacionar TierraViva con agricultura sostenible, eficiencia hídrica, consumo responsable y adaptación climática. \\
\hline
PCB & \textit{Printed Circuit Board} & Placa de circuito impreso planteada como posible mejora futura para aumentar robustez, orden eléctrico y facilidad de montaje. \\
\hline
PIN & \textit{Personal Identification Number} & Código de identificación asociado a la tarjeta SIM. En TierraViva se comprueba que la SIM permita el registro del módem LTE en la red móvil. \\
\hline
PK & \textit{Primary Key} & Clave primaria de una tabla de base de datos. En TierraViva se utiliza en el modelo SQLite para identificar de forma única registros de tablas como \texttt{readings}, \texttt{valve\_events} y \texttt{system\_state}. \\
\hline
PWM & \textit{Pulse Width Modulation} & Técnica de modulación utilizada para generar señales de control mediante variación del ciclo de trabajo. \\
\hline
RAM & \textit{Random Access Memory} & Memoria principal volátil utilizada durante la ejecución de un programa. En TierraViva se menciona al explicar por qué el estado de ingesta se conserva en SQLite y no únicamente en variables temporales. \\
\hline
RXD & \textit{Receive Data} & Línea o pin de recepción de datos del módulo A7670E-LASA en la comunicación serie con Arduino. \\
\hline
REST & \textit{Representational State Transfer} & Estilo de interfaz usado por la API HTTP local. \\
\hline
RF & Requisito funcional & Identificador de requisitos que describen funciones del sistema. \\
\hline
RFC & \textit{Request for Comments} & Serie de documentos técnicos utilizados para especificar estándares y protocolos de Internet. \\
\hline
RNF & Requisito no funcional & Identificador de requisitos de calidad, robustez, coste o reproducibilidad. \\
\hline
RSSI & \textit{Received Signal Strength Indicator} & Indicador de intensidad de señal móvil calculado a partir del valor CSQ. \\
\hline
RX & Recepción & Línea de recepción en comunicación serie. \\
\hline
SBC & \textit{Single-Board Computer} & Computador de placa única, como Raspberry Pi, utilizado en funciones de supervisión, almacenamiento, API o visualización. \\
\hline
SCL & \textit{Serial Clock} & Línea de reloj del bus I2C. \\
\hline
SDA & \textit{Serial Data} & Línea de datos del bus I2C. \\
\hline
SIM & \textit{Subscriber Identity Module} & Tarjeta usada por el módulo LTE para acceder a la red móvil. \\
\hline
SPI & \textit{Serial Peripheral Interface} & Bus serie mencionado en el contexto general de sistemas embebidos. \\
\hline
SQL & \textit{Structured Query Language} & Lenguaje utilizado para consultar la base de datos SQLite. \\
\hline
SRAM & \textit{Static Random Access Memory} & Memoria volátil disponible en Arduino UNO. Su capacidad limitada condiciona el tamaño de búferes, cadenas y mensajes construidos por el \textit{firmware}. \\
\hline
TFG & Trabajo de Fin de Grado & Marco académico del proyecto TierraViva. \\
\hline
TX & Transmisión & Línea de transmisión en comunicación serie. \\
\hline
TXD & \textit{Transmit Data} & Línea o pin de transmisión de datos del módulo A7670E-LASA en la comunicación serie con Arduino. \\
\hline
UART & \textit{Universal Asynchronous Receiver-Transmitter} & Comunicación serie usada conceptualmente entre microcontrolador y módem. \\
\hline
UDP & \textit{User Datagram Protocol} & Protocolo de transporte asociado a CoAP en el estado del arte. \\
\hline
ITU-T & \textit{International Telecommunication Union -- Telecommunication Standardization Sector} & Organismo de normalización citado por la recomendación Y.2060 para definir el concepto de IoT. \\
\hline
UNESCO & Organización de las Naciones Unidas para la Educación, la Ciencia y la Cultura & Sistema de clasificación científica empleado en el anexo de códigos UNESCO. \\
\hline
URL & \textit{Uniform Resource Locator} & Dirección usada para acceder a ThingSpeak, API o \textit{dashboard}. \\
\hline
USB & \textit{Universal Serial Bus} & Interfaz utilizada para programar Arduino o conectar dispositivos. \\
\hline
WAL & \textit{Write-Ahead Logging} & Modo de SQLite activado para mejorar la gestión de escrituras en la base de datos local. \\
\hline
WiFi & \textit{Wireless Fidelity} & Red inalámbrica local considerada como alternativa; en la arquitectura final de las estaciones remotas se prioriza LTE para la publicación de telemetría desde campo. \\
\hline
4G & Cuarta generación de comunicaciones móviles & Familia de tecnologías móviles dentro de la que se encuadra LTE, utilizada por el módulo A7670E-LASA para la publicación de telemetría. \\
\hline
\caption{Acrónimos empleados en la memoria de TierraViva.}
\label{tab:acronimos_tierraviva}\\
\end{longtable}

\normalsize

\section{Términos específicos de TierraViva}

Además de los términos generales anteriores, en la memoria se utilizan algunos nombres de variables, campos y estados propios del proyecto.

\renewcommand{\arraystretch}{1.50}
\scriptsize

\begin{longtable}{|p{0.30\textwidth}|p{0.56\textwidth}|}
\hline
\rowcolor{URJCRed}
\TVHead{Término} & \TVHead{Descripción} \\
\hline
\endfirsthead

\hline
\rowcolor{URJCRed}
\TVHead{Término} & \TVHead{Descripción} \\
\hline
\endhead

\hline
\multicolumn{2}{r}{\small Continúa en la página siguiente} \\
\endfoot

\hline
\endlastfoot

Estación A & Primera estación remota implementada y montada del prototipo. Publica datos en su canal ThingSpeak correspondiente y sirve como estación principal de validación del sistema. \\
\hline
Estación B & Segunda estación remota implementada y montada siguiendo la arquitectura base de TierraViva. Dispone de identificador y canal ThingSpeak propios, lo que permite separar sus datos de los de la estación A y validar el funcionamiento multiestación del sistema. \\
\hline
\texttt{main.py} & Proceso de ingesta ejecutado en Raspberry Pi. Consulta ThingSpeak, normaliza lecturas, evita duplicados y almacena nuevas entradas en SQLite. \\
\hline
\texttt{database.py} & Módulo encargado de crear y consultar la base de datos SQLite, gestionar tablas, índices, estado interno y operaciones de persistencia. \\
\hline
\texttt{api.py} & Módulo FastAPI que expone los \textit{endpoints} HTTP utilizados por el \textit{dashboard} para consultar datos almacenados en la Raspberry Pi. \\
\hline
\texttt{tierraviva.sh} & Script de gestión que automatiza instalación, arranque, parada, reinicio, comprobación de estado y consulta de logs del sistema software. \\
\hline
\texttt{tierraviva.db} & Fichero SQLite donde la Raspberry Pi almacena lecturas, eventos y estado interno del sistema de supervisión. \\
\hline
\texttt{field1} & Campo ThingSpeak asociado a \texttt{soil\_moisture\_pct}. \\
\hline
\texttt{field2} & Campo ThingSpeak asociado a \texttt{temperature\_c}. \\
\hline
\texttt{field3} & Campo ThingSpeak asociado a \texttt{humidity\_pct}. \\
\hline
\texttt{field4} & Campo ThingSpeak asociado a \texttt{light\_lux}. \\
\hline
\texttt{field5} & Campo ThingSpeak asociado a \texttt{pressure\_hpa}. \\
\hline
\texttt{field6} & Campo ThingSpeak asociado a \texttt{relay\_state}. \\
\hline
\texttt{field7} & Campo ThingSpeak asociado a \texttt{system\_event}. \\
\hline
\texttt{field8} & Campo ThingSpeak asociado a \texttt{debug\_code}. \\
\hline
\texttt{entry\_id} & Identificador de entrada asignado por ThingSpeak y utilizado para evitar duplicados. \\
\hline
\texttt{created\_at} & Fecha original de creación de una lectura en ThingSpeak. \\
\hline
\texttt{ts} & Marca temporal almacenada localmente en SQLite. \\
\hline
\texttt{station\_id} & Identificador de estación usado por la base de datos y la API. \\
\hline
\texttt{soil\_moisture\_pct} & Humedad del suelo expresada en porcentaje tras conversión desde lectura analógica. \\
\hline
\texttt{temperature\_c} & Temperatura ambiente en grados Celsius. \\
\hline
\texttt{humidity\_pct} & Humedad relativa ambiental. \\
\hline
\texttt{light\_lux} & Luminosidad medida en lux. \\
\hline
\texttt{pressure\_hpa} & Presión atmosférica medida en hPa. \\
\hline
\texttt{relay\_state} & Estado del relé publicado por la estación: apagado o activo. \\
\hline
\texttt{system\_event} & Código de evento que informa sobre reposo, riego, error, modo seguro o \textit{cooldown}. \\
\hline
\texttt{debug\_code} & Código auxiliar de diagnóstico publicado por la estación. \\
\hline
\texttt{raw\_json} & Respuesta original de ThingSpeak conservada para trazabilidad. \\
\hline
\texttt{source} & Origen de la lectura almacenada, normalmente \texttt{thingspeak}. \\
\hline
\texttt{readings} & Tabla principal de SQLite para almacenar lecturas. \\
\hline
\texttt{system\_state} & Tabla de SQLite usada para conservar estado de ingesta y últimos identificadores procesados. \\
\hline
\texttt{valve\_events} & Tabla para registrar eventos de actuación de riego. \\
\hline
\texttt{TIERRAVIVA\_HOME} & Ruta base del proyecto en la Raspberry Pi. \\
\hline
\texttt{STATIONS} & Lista de estaciones activas configuradas en el sistema. \\
\hline
\texttt{STATION\_ID} & Estación seleccionada por defecto por la API. \\
\hline
\texttt{THINGSPEAK\_CHANNEL\_A/B} & Identificador del canal ThingSpeak asociado a cada estación. \\
\hline
\texttt{THINGSPEAK\_READ\_KEY\_A/B} & Clave de lectura del canal ThingSpeak de cada estación. \\
\hline
\texttt{THINGSPEAK\_API\_KEY} & Clave de escritura utilizada por el \textit{firmware} de estación. Debe mantenerse fuera del repositorio público. \\
\hline
\texttt{POLL\_INTERVAL\_SECONDS} & Intervalo de consulta de ThingSpeak por parte de la Raspberry Pi. \\
\hline
\texttt{REQUEST\_TIMEOUT\_SECONDS} & Tiempo máximo de espera en peticiones HTTP. \\
\hline
\texttt{SOIL\_LOW\_THRESHOLD} & Umbral inferior de humedad usado para clasificar el suelo como seco. \\
\hline
\texttt{SOIL\_HIGH\_THRESHOLD} & Umbral superior de humedad usado para clasificar el suelo como húmedo. \\
\hline
\texttt{WEB\_REFRESH\_SECONDS} & Intervalo recomendado de actualización del \textit{dashboard}. \\
\hline
\texttt{DATA\_STALE\_SECONDS} & Tiempo usado para decidir cuándo una lectura deja de considerarse reciente. \\
\hline
\caption{Términos específicos utilizados en TierraViva.}
\label{tab:terminos_especificos_tierraviva}
\end{longtable}

\normalsize

\section{Códigos de estado y diagnóstico}

Los códigos publicados por la estación permiten interpretar el estado operativo sin depender únicamente de las variables ambientales.

\begin{table}[H]
\centering
\renewcommand{\arraystretch}{1.25}
\small
\begin{tabular}{|p{0.16\textwidth}|p{0.28\textwidth}|p{0.42\textwidth}|}
\hline
\rowcolor{URJCRed}
\TVHead{Código} & \TVHead{Campo habitual} & \TVHead{Interpretación} \\
\hline
\texttt{204} & \texttt{debug\_code} & No se ha activado el riego en el ciclo actual. \\
\hline
\texttt{200} & \texttt{debug\_code} & Funcionamiento correcto. \\
\hline
\texttt{204} & \texttt{debug\_code} & No se ha activado el riego en el ciclo actual. \\
\hline
\texttt{206} & \texttt{system\_event} / \texttt{debug\_code} & Modo seguro activo. \\
\hline
\texttt{210} & \texttt{system\_event} & Riego activado. \\
\hline
\texttt{211} & \texttt{system\_event} & Riego finalizado y relé apagado. \\
\hline
\texttt{400} & \texttt{system\_event} / \texttt{debug\_code} & Error relevante, normalmente asociado a sensores o lectura no válida. \\
\hline
\texttt{429} & \texttt{system\_event} / \texttt{debug\_code} & Riego bloqueado por \textit{cooldown}. \\
\hline
\end{tabular}
\caption{Códigos de estado y diagnóstico utilizados por TierraViva.}
\label{tab:codigos_estado_diagnostico_glosario}
\end{table}
